\section{参考资料：课程给顺序，教材给边界，讲义给解释}

参考资料要为正文承担不同角色，单纯堆积链接没有意义。大学课程最适合观察教学顺序和典型问题；教材与专著负责定义、定理条件和完整证明；公开讲义常能补足直觉、推导与例子；论坛只适合发现读者反复困惑的地方，不能单独证明数学结论。

当来源表述不一致时，应先检查定义域、量词、正则条件和记号，再查作者或出版社勘误。仍无法判断，就回到原始论文或亲自完成推导，而不是以网页数量投票。

\subsection{各学科的主要入口}

{\def\LTcaptype{none} % do not increment counter
\small
\begin{longtable}{@{}>{\raggedright\arraybackslash}p{2.2cm}>{\raggedright\arraybackslash}p{4cm}>{\raggedright\arraybackslash}p{4cm}>{\raggedright\arraybackslash}p{4cm}@{}}
\toprule\noalign{}
学科 & 课程主线 & 严格参考 & 补充讲解 \\
\midrule\noalign{}
\endhead
\bottomrule\noalign{}
\endlastfoot
数学语言 & MIT 6.042J Mathematics for Computer Science & Discrete Mathematics: An Open Introduction & MIT 6.042J Readings \\
微积分 & MIT 18.01SC 与 MIT 18.02SC & OpenStax Calculus & Paul\textquotesingle s Online Math Notes \\
线性代数 & MIT 18.06 Linear Algebra & Linear Algebra --- Jim Hefferon & MIT 18.06 Video Lectures \\
矩阵函数与张量 & MIT 18.06 与 MIT 18.065 & Functions of Matrices: Theory and Computation --- Nicholas J. Higham & Tensor Decompositions and Applications --- Tamara G. Kolda \& Brett W. Bader \\
离散数学 & MIT 6.042J 与 Berkeley CS 70 & Discrete Mathematics: An Open Introduction & Mathematics for Computer Science \\
群、布尔代数与自动机 & MIT 18.701 与 MIT 6.045J & Algebra --- Michael Artin；Introduction to the Theory of Computation --- Michael Sipser & Visual Group Theory --- Nathan Carter \\
概率论 & Berkeley EECS 126 & Introduction to Probability --- Bertsekas \& Tsitsiklis & Harvard Stat 110 \\
随机过程 & MIT 6.041SC Unit III & Introduction to Probability, Statistics, and Random Processes & MIT RES.6-012 Part III \\
随机微分方程与极值 & MIT 18.676 Stochastic Calculus & Stochastic Differential Equations: An Introduction with Applications --- Bernt Øksendal & Modelling Extremal Events for Insurance and Finance --- Embrechts, Klüppelberg \& Mikosch \\
信息论 & Stanford EE 276 & Information Theory, Inference, and Learning Algorithms --- MacKay & MIT 6.441 Lecture Notes \\
数值分析 & MIT 18.330 & Fundamentals of Numerical Computation & Nick Trefethen\textquotesingle s Numerical Analysis Lectures \\
时频分析与小波 & Stanford EE 261 & A Wavelet Tour of Signal Processing: The Sparse Way --- Stéphane Mallat & Ten Lectures on Wavelets --- Ingrid Daubechies \\
凸优化 & Stanford EE364A & Convex Optimization --- Boyd \& Vandenberghe & EE364A Lecture Materials \\
矩阵微积分 & Stanford CS229 & Matrix Differential Calculus --- Magnus \& Neudecker & The Matrix Calculus You Need for Deep Learning \\
数理统计 & Stanford STATS 200 & All of Statistics --- Larry Wasserman & Penn State STAT 415 \\
统计决策与可信预测 & Stanford STATS 300A & Statistical Decision Theory and Bayesian Analysis --- James O. Berger & Algorithmic Learning in a Random World --- Vovk, Gammerman \& Shafer \\
实分析支撑 & MIT 18.100B Real Analysis & Understanding Analysis --- Stephen Abbott & UC Davis Analysis Notes --- John Hunter \\
机器学习 & Stanford CS229 与 Berkeley CS189 & Probabilistic Machine Learning --- Kevin Murphy & Stanford Engineering Everywhere CS229 \\
因果推断 & HarvardX Causal Diagrams & Causal Inference: What If --- Hern\'an \& Robins & Causality: Models, Reasoning, and Inference --- Judea Pearl \\
深度学习 & Stanford CS231n 与 CS224N & Deep Learning --- Goodfellow, Bengio \& Courville & Dive into Deep Learning \\
统计学习与核方法 & Stanford CS229T & Foundations of Machine Learning --- Mohri, Rostamizadeh \& Talwalkar & Learning with Kernels --- Schölkopf \& Smola \\
微分方程与动力系统 & MIT 18.03SC Differential Equations & Nonlinear Dynamics and Chaos --- Steven H. Strogatz & Differential Equations, Dynamical Systems, and an Introduction to Chaos --- Hirsch, Smale \& Devaney \\
最优控制与强化学习 & Stanford EE363 与 Stanford CS234 & Dynamic Programming and Optimal Control --- Dimitri P. Bertsekas & Reinforcement Learning: An Introduction --- Sutton \& Barto \\
\end{longtable}
}

链接优先指向大学、作者或出版社入口，不在知识库中镜像受版权保护的教材。引用只保留理解所需的短句，其余用自己的语言重新推导和解释。公开可访问不等于可以大段复制。

论坛与问答可以帮助发现``为什么条件不能省''\,``两个术语有何区别''之类问题。得到线索后仍要回到课程、教材或原始论文；赞成票和被接受回答都不是数学证据。
